\chapter{Confluence, Critical Pairs, and Completion}
\label{ch:rewrite-confluence}

Running the same simplifier with a different worklist order should not unexpectedly change its canonical answer.
Termination prevents infinite search but does not prevent two different normal forms.
Confluence addresses this independence from reduction choice.
Critical pairs reduce an infinite family of possible conflicts to finitely described overlaps for a finite first-order rule collection, under precisely stated conditions.

\section{Analyze a local peak}
\label{sec:ch84-peak}
Take rules $f(a)\to b$ and $a\to c$.
At term $f(a)$, a root rewrite produces b and a child rewrite produces $f(c)$.
If there are no further rules, both results are distinct normal forms.
Figure~\ref{fig:ch84-peak} draws the failed joining obligation.
Neither changing the rule declaration order nor always rewriting children first makes the full relation confluent.

\begin{figure}[htbp]
\centering
\begin{tikzpicture}[node distance=15mm and 17mm]
\node[basic block] (a) {f(a)};
\node[basic block,below left=of a] (b) {b: normal};
\node[basic block,below right=of a] (c) {f(c): normal};
\draw[cfg edge] (a)--node[left,font=\small]{root}(b);
\draw[cfg edge] (a)--node[right,font=\small]{child}(c);
\end{tikzpicture}
\caption{A root/child overlap creates distinct normal forms unless further rules join them.}
\label{fig:ch84-peak}
\end{figure}

Two redexes at disjoint tree positions commute because replacing one leaves the other's subterm intact.
Nested redexes demand more analysis.
An overlap inside a left-side variable is handled by substitution reasoning in the ordinary critical-pair theorem; overlaps at nonvariable positions produce the explicit critical pairs.
Repeated variables and nonlinear patterns require the theorem's actual formulation, not a blanket assertion that every nested replacement commutes.

\section{Rename variables before unifying overlaps}
\label{sec:ch84-enumeration}
For rules $l_1\to r_1$ and $l_2\to r_2$, first rename variables apart.
At every nonvariable position $p$ of $l_2$, unify $l_1$ with $l_2|_p$.
For a most-general unifier $\sigma$, the critical pair is
\[
 \bigl(l_2[r_1]_p\sigma,\;r_2\sigma\bigr).
\]
The two terms represent the nested-rule step and the outer-rule step from their common overlap instance.
Variable renaming is essential: identical metavariable spellings in independently declared rules do not impose a shared binding before overlap construction.

Algorithm~\ref{alg:ch84-pairs} is implemented in the companion.
It omits trivial equal pairs but retains nontrivial duplicates for transparency.
For $R$ rules with at most $N$ left-side positions each, it attempts $O(R^2N)$ unifications, with actual unification and instantiated-output costs added.
The elementary unifier's costs are not hidden inside a false constant-time bound.

\begin{algorithm}[htbp]
\caption{Enumerate first-order critical overlaps with disjoint variable namespaces.}
\label{alg:ch84-pairs}
\begin{pseudocode}
\Require finite well-formed first-order rewrite rules
\Ensure critical pairs from all nonvariable left-side overlaps
\ForAll{ordered rule pairs $(l_1\to r_1,l_2\to r_2)$}
\State rename the two variable sets apart
\ForAll{nonvariable positions $p$ in $l_2$}
\State attempt finite-term unification of $l_1$ and $l_2|_p$
\If{unifier $\sigma$ exists}
\State emit $(l_2[r_1]_p\sigma,r_2\sigma)$ unless both terms are equal
\EndIf
\EndFor
\EndFor
\end{pseudocode}
\end{algorithm}

For the worked rules, unifying a with the child a of f(a) yields the empty substitution.
The two branches are f(c) and b.
Table~\ref{tab:ch84-pairs} also contrasts the add-zero/multiply-one system, whose nontrivial overlap list is empty.
An empty list in that particular ordinary first-order system establishes local confluence through the critical-pair theorem, not arbitrary semantic correctness of the identities.

\begin{table}[htbp]
\caption{Critical overlaps identify local rule interactions rather than search-order preferences.}
\label{tab:ch84-pairs}
\centering
\begin{tabular}{lll}
\toprule
Rule set & Nontrivial overlap & Joinability here \\
\midrule
f(a) to b; a to c & b versus f(c) & no further rule; not joined \\
add-zero; multiply-one & none & local confluence follows \\
first set plus f(c) to b & b versus f(c) & joined at b \\
\bottomrule
\end{tabular}
\end{table}

\section{Use termination when reducing both branches}
\label{sec:ch84-join}
For a terminating system, every reduction reaches some normal form.
If the system is already confluent, reducing both sides of a pair and comparing normal forms decides their joinability.
When confluence is the property being checked, arbitrary chosen normal forms that differ do not by themselves prove nonjoinability: different paths might reach a common descendant before or after those choices in a nonconfluent system.
One must use a complete joinability procedure for the setting or provide actual joining reductions.

The common sufficient workflow is to exhibit a joining reduction for every critical pair and separately establish termination.
The critical-pair theorem yields local confluence, and the terminating/local-confluent theorem yields confluence~\cite{ohlebusch-2002-rewriting}.
For a terminating finitely branching relation, exhaustive descendant exploration can decide a finite pair's joinability, but its search space can be very large.
Bounded failed search is unresolved, not a proof that no join exists.

An \term{orthogonal} first-order system is left-linear and has no nontrivial overlaps in the appropriate definition.
Orthogonality yields confluence without a termination requirement.
This makes it attractive for functional evaluation rules, but normal-form existence and strategy behavior remain separate questions.
Adding an optimization rule can destroy nonoverlap, so an orthogonality claim must refer to the complete combined rule set.

\section{Complete equations into oriented rules}
\label{sec:ch84-completion}
\term{Completion} attempts to turn equations into a terminating confluent presentation by orienting them under a reduction ordering, simplifying them, and adding equations for unresolved critical pairs.
Figure~\ref{fig:ch84-loop} shows the feedback loop.
The ordering is an input to the proof discipline, not an arbitrary preference for shorter printed text.
New rules can create new overlaps and invalidate old simplified forms, which must be revisited.

\begin{figure}[htbp]
\centering
\begin{tikzpicture}[node distance=13mm]
\node[basic block] (e) {pending equations};
\node[basic block,below=of e] (r) {simplify and orient};
\node[basic block,below=of r] (p) {new rules and critical pairs};
\draw[cfg edge] (e)--(r);\draw[cfg edge] (r)--(p);
\draw[cfg edge] (p.east) to[bend right=50] node[right,font=\small]{new obligations}(e.east);
\end{tikzpicture}
\caption{Completion creates new obligations whenever it adds or changes a rule.}
\label{fig:ch84-loop}
\end{figure}

Algorithm~\ref{alg:ch84-completion} is a conceptual fair completion loop, not an implemented complete solver in the lab.
It records orientation failure and budget exhaustion explicitly.
Success additionally requires that all generated obligations were processed fairly and that the maintained rules/equations preserve the original equational theory under the chosen completion inference system.
Merely obtaining an empty queue after accidentally dropping a pair is not a theorem.

\begin{algorithm}[htbp]
\caption{Attempt completion while retaining orientation and overlap obligations.}
\label{alg:ch84-completion}
\begin{pseudocode}
\Require equations, reduction ordering, sound completion inferences, resource budget
\Ensure certified convergent presentation on success, or explicit unresolved outcome
\While{an obligation remains and resources permit}
\State fairly select and simplify an equation under current rules
\If{its sides are identical} \State discharge it
\ElsIf{one side is greater under the reduction ordering}
\State orient; interreduce affected rules; enqueue all new required overlaps
\Else \State retain unorientable obligation or return unresolved
\EndIf
\EndWhile
\State claim success only if every obligation is discharged and certificates check
\end{pseudocode}
\end{algorithm}

Completion can fail to orient an equation or run indefinitely while generating larger obligations.
There is no finite general upper bound promising successful completion of arbitrary input theories.
For the simple peak, adding f(c)$\to$b joins the conflict.
The weighted interpretation $a=2$, $c=1$, $b=0$, and $f(x)=x+1$ strictly decreases all three rules; ordinary node count alone does not decrease a$\to$c.
That worked repair illustrates the idea without claiming that the lab contains a general completion implementation.

\section{Apply canonicalization claims carefully in compilers}
\label{sec:ch84-compilers}
A compiler often needs only a deterministic profitable result, not a unique normal form for every semantically equivalent program.
Demanding full confluence of every optimization rule can be expensive and unnecessary when phase ordering is deliberate.
Conversely, a canonical IR used for equality hashing or persistent caches needs a much clearer uniqueness contract.
Conditional rules involving analysis facts complicate overlaps because their premises may not hold simultaneously or may become stale after rewriting.

Table~\ref{tab:ch84-contracts} separates the possible promises.
Equality saturation keeps alternatives instead of demanding destructive rewrite confluence, but still needs sound equalities and a finite extraction policy.
The earlier e-graph chapter is therefore complementary, not a replacement for critical-pair reasoning about a normalizer.
The optimizer's tests should state whether they check semantic agreement, deterministic output, idempotence, or true canonicality.

\begin{table}[htbp]
\caption{Canonicalization is a stronger promise than deterministic optimization.}
\label{tab:ch84-contracts}
\centering
\begin{tabular}{ll}
\toprule
Promise & Required argument \\
\midrule
semantic preservation & every applied transformation is sound \\
deterministic result & fixed rule, match, and worklist ordering \\
idempotence & rerunning the selected pass changes nothing \\
unique normal form & confluence for elements that normalize \\
convergent equality decision & termination, confluence, theory correspondence \\
\bottomrule
\end{tabular}
\end{table}

\section{See how variable-position peaks are joined}
\label{sec:ch84-variable-peaks}
Nonvariable overlap enumeration does not mean reductions inside substituted variables can be ignored without an argument.
Consider f(x,x)$\to$c and a$\to$b at source f(a,a).
The root step gives c.
Reducing only the first a gives f(b,a), where the nonlinear root pattern no longer matches.
Reducing the remaining a yields f(b,b), which again matches the root rule and reduces to c.
The joining sequence must synchronize the affected repeated occurrences; simply claiming that one child step commutes with the root step would be inaccurate.

Duplication on the right has the dual effect.
For f(x)$\to$g(x,x), an inner step in the substituted x may need to be repeated in both copies after the root step.
This is why a graph update shared by two incoming edges can correspond to several tree steps rather than exactly one.
The ordinary critical-pair theorem packages these variable-overlap arguments together with the explicitly enumerated nonvariable peaks.
Its proof assumptions must match whether the engine works on trees, shared graphs, conditional rules, or rewriting modulo equations.

For a compiler, side conditions can make critical-pair construction require more than unifying syntax.
Two rules may overlap structurally but have mutually inconsistent range conditions.
Alternatively, applying one can destroy the analysis premise of the other.
A correct conditional analysis retains those conditions as obligations rather than dropping them or assuming every structural pair is jointly feasible.
The finite unconditional lab remains useful precisely because it keeps that extra solver layer outside its stated theorem.

\section{Treat completion output as a certificate-bearing artifact}
\label{sec:ch84-certificates}
A successful completion run should retain why each final rule follows from the original equations and why each orientation decreases the supplied ordering.
It should also retain the joining reductions for critical pairs under the final rule set or under a sound incremental inference discipline.
A small checker can then verify these local facts without replaying every heuristic choice used by the completion search.
Search sophistication and certificate-checker complexity need not grow together.

Interreduction complicates bookkeeping.
If a rule's right side is simplified, its proof changes by a known conversion.
If its left side becomes reducible by another rule, keeping both old and new versions without a sound replacement step can create redundant or misleading overlap obligations.
The inference system must preserve the original equational theory while updating the active rule set.
A transcript containing only the final printed rules is insufficient to reconstruct those reasons unless the checker independently solves the full equational problem.

Completion is therefore a good example of the book's general trust split.
Heuristics choose equations, orientable alternatives, simplification opportunities, and resource allocation.
The checker establishes the admitted mathematical claims.
Budget exhaustion can still produce useful proved rules, but it cannot be published as a convergent presentation when unresolved overlaps remain.
Likewise, a compiler may use a sound incomplete simplifier without claiming a complete canonical equality decision for its source language.

\section{Exercises}
\label{sec:ch84-exercises}
\begin{enumerate}
\item\label{ex:ch84-peak} Construct the critical pair for f(a)$\to$b and a$\to$c, including its overlap position.
\item\label{ex:ch84-repair} Add f(c)$\to$b and give explicit joining reductions and a weighted-interpretation termination argument.
\item\label{ex:ch84-apart} Explain why variables in two copies of the same rule must be renamed apart before computing self-overlaps.
\item\label{ex:ch84-budget} A critical-pair join search runs out of fuel with unequal current terms.
What conclusion is justified, and what additional evidence would prove joinability?
\end{enumerate}

\section{Further reading}
\label{sec:ch84-reading}
Ohlebusch develops critical-pair criteria, confluence, and completion beyond the finite examples used here~\cite{ohlebusch-2002-rewriting}.
Dershowitz provides the ordering framework that prevents oriented completion rules from introducing infinite descent~\cite{dershowitz-1982-orderings}.
